\begin{tikzpicture}[x=10mm,y=10mm,line width=.65pt,font=\fontsize{8}{10}\selectfont]
\path[use as bounding box] (0,-.35) rectangle (10,3.6);
\draw[->,Line,line width=.8pt] (.6,2.4)--(8.6,2.4);
\node at (9.35,2.4) {迭代次数$n$};
\draw[->,Line,line width=.8pt] (.6,.35)--(.6,3.05);
\node at (2.55,3.25) {误差$e_n=d(x^{(n)},x_*)$};
% 指数收缩误差曲线（折线示意）
\draw[Blue,line width=1.1pt] plot[smooth] coordinates {(1,2.8) (1.7,2.45) (2.4,2.15) (3.1,1.9) (3.8,1.68) (4.5,1.5) (5.2,1.35) (5.9,1.23) (6.6,1.13) (7.3,1.06) (8,1.01) (8.7,.98)};
% 迭代点
\foreach \p in {(1,2.8),(2.4,2.15),(3.8,1.68),(5.2,1.35),(6.6,1.13),(8,.98)} {\fill[Ink] \p circle[radius=1mm];}
\node[Blue] at (5.6,2.9) {$e_n\leq q^n(1-q)^{-1}d(x^{(1)},x^{(0)})$};
% q 提示
\node[Muted] at (7.6,.5) {$q<1$，误差几何收缩到$0$};
\end{tikzpicture}
